% =====================================================================
%  29.4  Cactus graphs: cycles that never share an edge.
%  Exports: sec:c31-cactus  def:c31-cactus  thm:c31-cactustree
%  lem:c31-cycleorder  fig:c31-cactus  lst:c31-cactusdp
%  WORKED section + a solved lab.  A cactus is an undirected graph in which every
%  edge lies in at most one simple cycle (equivalently every block is a bridge or
%  a simple cycle).  Contracting each cycle-block to a node gives a tree (the
%  block-cut tree specialised), so DP over a cactus = DP over that tree with one
%  extra primitive: a "cycle transfer" that folds a cycle node's children around
%  the ring.  Running object: the bowtie (two triangles sharing one vertex),
%  counted for independent sets -> 10.  In-section solved lab: CF 231E "Cactus"
%  (2100, L5) -- a VERTEX cactus (each vertex on at most one cycle, so cycles are
%  vertex-disjoint); for k queries (x,y) count simple x-y paths, which is 2^c with
%  c the number of cycles on the x-y route.  Solved by DSU-contracting each cycle,
%  building the contracted tree, and answering c by LCA + a root-to-node cycle
%  prefix.  Live-verified against the CF page (2s, 256 MB, n,m,k<=1e5, mod 1e9+7);
%  the single official sample (10-vertex graph, 6 queries) is replayed in full
%  -> 2 2 2 4 4 1.  Promoted from this chapter's mixed set, so its URL now lives
%  here ONLY (no duplicate).
% =====================================================================
\section{Cactus Graphs: Cycles That Never Share an Edge}\label{sec:c31-cactus}

\begin{previouslybox}[Blocks were general; here every block is a cycle or an edge]
\(\secref{c31-block}\) decomposed any graph into blocks via the edge stack.  A
\term{cactus} is the family where those blocks are as simple as possible: each is
either a single bridge or a single \emph{simple cycle} (no block is a tangled
\(K_4\)).  That restriction means every edge sits in at most one cycle, so you can
contract each whole cycle to one node and still have a \emph{tree}
(\(\thmref{c31-cactustree}\)).  Dynamic programs that are hard on a general graph --
independent set, dominating set, matching, colourings -- are polynomial on a cactus
because they are a tree DP plus one local cycle computation
(\(\lemref{c31-cycleorder}\)).
\end{previouslybox}

\begin{topicintro}
Formally a \term{cactus} is a connected graph where any two simple cycles share at
most one vertex (\(\defref{c31-cactus}\)); equivalently no edge lies in two cycles.
The block-cut tree of \(\secref{c31-block}\) specialises: block-nodes are bridges and
cycles, cut-vertex-nodes are the articulation points where cycles touch.  Collapse
each cycle to a single node and you get an ordinary tree, and a bottom-up DP is legal
on it -- the only new operation is, at a cycle block, combining the in-tree values
that hang around the ring \emph{in cyclic order} with one endpoint's state pinned by
the parent.  That cycle transfer (\(\lemref{c31-cycleorder}\)) is the entire extra
craft.

\textbf{Prerequisites.} \(\secref{c31-block}\) (blocks, the edge stack); Ch.~13
(\chref{c10}) (the fold); the transfer is a bounded ring DP, no new machinery.

\textbf{After this section you can}
\begin{itemize}[nosep]
  \item test the cactus property while popping blocks (reject a block that is not a cycle/bridge); \hfill\lv{4}
  \item contract each cycle block and recognise the resulting tree; \hfill\lv{4}
  \item write the cycle-transfer DP step that folds a ring with one pinned state; \hfill\lv{5}
  \item count independent sets / colourings of a cactus in \(\bigO{n}\). \hfill\lv{4}
\end{itemize}
\end{topicintro}

% ---------------------------------------------------------------
\subsection{Cycles as blocks}
% ---------------------------------------------------------------

\begin{definition}{Cactus graph}{c31-cactus}
A connected undirected graph is a \term{cactus} if every edge belongs to at most one
simple cycle, equivalently if every block of its block-cut tree
(\(\defref{c31-cutv}\)) is either a single bridge or a single simple cycle.  Any two
cycles then share at most one vertex.
\end{definition}

\begin{theorem}{The contracted cactus is a tree}{c31-cactustree}
Given a cactus, replace each maximal simple cycle by one node (keeping the
attachment vertices' incidences) and keep bridges as edges.  The result is a tree.
\end{theorem}

\begin{proofbox}
Blocks of any graph already form a tree via cut vertices (\(\defref{c31-bct}\)); the
cactus condition only restricts \emph{which} blocks appear (cycles and bridges), not
how they glue.  Two cycles sharing an edge would form one non-cycle block, forbidden,
so distinct cycles share at most a cut vertex, which the block-cut tree already
models as a single connecting node -- no cycle of block-nodes can close.  Contracting
each cycle-block to a node and each bridge to an edge yields a connected graph on
\(\le\) (blocks) nodes with (blocks\(-1\)) inter-block links, i.e.\ a tree. \(\Box\)
\end{proofbox}

\begin{lemma}{Cyclic-order transfer}{c31-cycleorder}
Let a cycle block \(C=(c_0,c_1,\dots,c_{k-1},c_0)\) be entered from its parent through
the articulation \(c_0\), with in-tree (already-computed) child values \(a_i\) at each
\(c_i\).  A DP with a per-vertex state set \(S\) and a legality/weight table
\(\mathit{mix}\) combines them by a linear ring DP that pins \(c_0\)'s state to the
value the parent fixed and folds \(c_1,\dots,c_{k-1}\) in order, closing only if the
last is consistent with the pinned \(c_0\).  The cost is \(\bigO{k\cdot|S|^2}\) per
cycle and the sum of \(k\) over cycles is \(\bigO n\).
\end{lemma}

\begin{lstlisting}[caption={Cactus DP skeleton: pop blocks; a bridge is a tree-edge update, a cycle runs a ring transfer with the entry vertex pinned to each candidate state and the best taken. States here are in/out for independent sets. ASCII only, identifiers without underscores}, label={lst:c31-cactusdp}]
// dp0[v], dp1[v]: # ways for the sub-cactus rooted at v with v OUT / IN.
// Bridge block u-v (v child):
void doBridge(int u,int v){
  dp0[u]*= dp0[v]+dp1[v];
  dp1[u]*= dp0[v];
}
// Cycle block with entry c0 (state pinned to s) and ring c1..c_{k-1}:
long long ring(int k, vector<long long>&z0, vector<long long>&z1, int s){
  // fold c1..c_{k-1}; 'prev0/prev1' carry the running pair, adjacency = no two IN
  long long prev0=1, prev1=0;              // c0 fixed: if s==IN, c1 forced OUT
  for(int i=1;i<k;i++){
    long long n0 = (prev0+prev1) * z0[i];
    long long n1 = prev0 * z1[i];          // c_i IN only if c_{i-1} OUT
    prev0=n0; prev1=n1;
  }
  if(s==1 && prev1>0) return 0;            // closing edge c_{k-1}-c0 violated
  return prev0+prev1;
}
// At a cycle node u (=c0) combine both pins and multiply into u's DP:
void doCycle(int k, vector<long long>&z0, vector<long long>&z1){
  long long out = ring(k,z0,z1,0);         // c0 OUT
  long long inW = ring(k,z0,z1,1);         // c0 IN
  // caller multiplies dp0/dp1 of u's parent by (out) and (inW) as appropriate
}
\end{lstlisting}

\begin{figure}[H]
\centering
\begin{tikzpicture}[diagram, scale=0.95, transform shape,
   cyc/.style={circle, draw=softpurple!80, fill=softpurple!12, line width=1pt,
               minimum size=6.6mm, font=\scriptsize\bfseries, inner sep=1pt},
   cut/.style={circle, draw=difforange, fill=difforange!20, line width=1.3pt,
               minimum size=6.6mm, font=\scriptsize\bfseries, inner sep=1pt}]
  \node[cut] (v) at (0,0) {v};
  \node[cyc] (x) at (-1.2,1.0) {x}; \node[cyc] (y) at (-1.2,-1.0) {y};
  \node[cyc] (p) at ( 1.2,1.0) {p}; \node[cyc] (q) at ( 1.2,-1.0) {q};
  \draw[softpurple!75, line width=1.1pt] (v)--(x)--(y)--(v);
  \draw[softpurple!75, line width=1.1pt] (v)--(p)--(q)--(v);
  \node[annot, font=\scriptsize] at (-1.2,1.9) {triangle 1};
  \node[annot, font=\scriptsize] at ( 1.2,1.9) {triangle 2};
\end{tikzpicture}
\caption{The \term{bowtie} cactus: two triangles \(\{v,x,y\}\) and \(\{v,p,q\}\) sharing a
single cut vertex \(v\) (orange).  Each edge lies in exactly one cycle, so it is a cactus
(\(\defref{c31-cactus}\)); contracting each triangle to a node leaves the two-node tree
joined at \(v\) (\(\thmref{c31-cactustree}\)).  Its independent-set count, \(10\,\), is the
worked trace below.}
\label{fig:c31-cactus}
\end{figure}

\subsection{Worked trace: independent sets of the bowtie}

Take the \term{bowtie}: two triangles sharing exactly one vertex \(v\), with
triangle~1 \(=\{v,x,y\}\) and triangle~2 \(=\{v,p,q\}\); every edge is in one cycle,
so it is a cactus (\(\defref{c31-cactus}\)).  Count its independent sets with the
contracted-tree DP.

\begin{tracebox}[Root at the shared vertex \(v\): the two cycle transfers multiply]
Each triangle is a cycle block entered at \(v\).  \textbf{Pin \(v\) IN}
(\(s{=}1\)): in triangle \(\{v,x,y\}\) both \(x,y\) are forced OUT, contributing
\(1\); likewise triangle \(\{v,p,q\}\) contributes \(1\).  So
\(\mathit{dp}_1(v)=1\) (the set \(\{v\}\)).  \textbf{Pin \(v\) OUT} (\(s{=}0\)): in
triangle 1 the ring over \(x,y\) with \(v\) out allows \(\emptyset,\{x\},\{y\}\)
-- \(3\) independent sets (the closing edge \(x\!-\!y\) forbids \(\{x,y\}\)); the
transfer \code{ring(3,\ldots,0)} returns \(3\).  Triangle 2 likewise \(3\).  So
\(\mathit{dp}_0(v)=3\cdot3=9\).  \textbf{Total independent sets} \(=\mathit{dp}_0(v)+
\mathit{dp}_1(v)=9+1=\boxed{10}\).  Cross-check by hand: either \(v\) is chosen
(then \(x,y,p,q\) all excluded: \(1\) set) or not (then each edge \(x\!-\!y\),
\(p\!-\!q\) independently offers \(\emptyset/\text{one endpoint}\): \(3\times3=9\)).
\(10\) confirmed. \(\Box\)
\end{tracebox}

\begin{insightbox}[Cactus versus the general block-cut tree]
On a general graph a block can be a dense 2-connected blob with exponentially many
local configurations, so a block-node DP is hopeless.  The cactus promise -- every
block is a cycle -- is exactly what bounds each block's local DP at
\(\bigO{k\cdot|S|^2}\): a cycle is a ring, and a ring is one linear pass.  This is
why problems phrased as ``the graph is a cactus'' or ``\(m\le n+O(1)\) and it is
connected enough to be one'' are tractable: the whole difficulty is pushed into
independent per-cycle transfers glued by Ch.~13's tree fold.
\end{insightbox}

\begin{complexitybox}[Cactus DP]
Detecting the cactus property is one edge-stack pass: pop blocks
(\(\thmref{c31-bccpop}\)) and \emph{reject} if any block has \(|E|\ne|V|\) and
\(|E|>1\) (a non-cycle, non-bridge block).  Then a postorder over the contracted
tree does one transfer per cycle: total \(\bigO{n\cdot|S|^2}\) (constant
\(|S|\) for independent set / 2-colouring, so linear in \(n\)).  Space \(\bigO n\).
\end{complexitybox}

\begin{warningbox}[Three cactus errors]
\begin{enumerate}[nosep,leftmargin=1.6em]
  \item \textbf{Treating a shared-edge pair of cycles as a cactus.} Two triangles
        sharing an \emph{edge} (a diamond) are \emph{not} a cactus -- the shared edge is
        on two cycles, so that block has \(|E|>|V|\); the transfer of
        \(\lemref{c31-cycleorder}\) is invalid there.
  \item \textbf{Forgetting the closing edge.} In \code{ring}, after folding
        \(c_1..c_{k-1}\) you must reject the configurations where \(c_{k-1}\) conflicts
        with the pinned \(c_0\); dropping it turns a cycle into a path and over-counts.
  \item \textbf{Double-counting the articulation.} The entry vertex \(c_0\) belongs to
        the cycle block \emph{and} to its parent block; pin its state once and reuse, or
        you multiply its contribution twice.
\end{enumerate}
\end{warningbox}

\begin{hardeningbox}
\begin{itemize}[nosep,leftmargin=1.4em]
  \item \textbf{Generalise the state, keep the shape.} The bowtie trace is independent
        set; swap \(\mathit{mix}\) for a proper-colouring table (states = colours) and
        \code{ring} counts colourings; a dominating-set table counts domination.  Only
        the constant \(|S|\) and the \code{prev0/prev1} recurrence change.
  \item \textbf{Near-cactus inputs.} If the graph is not a cactus but has few extra
        edges, condition on the endpoints of the offending edges (\(\bigO{2^{\text{few}}}\))
        to reduce to a cactus -- the classic ``cyclomatic number as a parameter'' trick.
  \item \textbf{Build the block tree explicitly.} For many cycles, materialise the
        contracted tree (block-node per cycle, cut-vertex nodes) and run a normal
        tree-DP driver, calling \code{doCycle} at cycle nodes -- you reuse Ch.~13's
        postorder instead of hand-crafting a cactus walk.
\end{itemize}
\end{hardeningbox}

% ---------------------------------------------------------------
\subsection{Solved: CF 231E -- Cactus (2100, L5)}
% ---------------------------------------------------------------

\solvedproblem{c31-cactuscount}{Cactus}{CF 231E}{https://codeforces.com/problemset/problem/231/E}{CF 2100 L5}{vertex cactus, DSU-contract each cycle, contracted tree, LCA + cycle-prefix, count routes as \(2^{c}\)}

\begin{problemstatement}[CF 231E, verified against the live CF problem page]
A connected undirected graph on \(n\) vertices and \(m\) edges is a \term{vertex
cactus}: every vertex belongs to at most one simple cycle (so distinct cycles are
\emph{vertex-disjoint}).  There are no loops or multiple edges.
A \term{simple path} is a walk that uses each edge at most once; two paths are
distinct when their \emph{edge sets} differ.  Given \(k\) queries \((x_i,y_i)\)
(\(x_i\neq y_i\)), for each count the number of distinct simple paths from \(x_i\)
to \(y_i\), modulo \(10^9+7\).  Constraints: \(2\le n\le 10^5\), \(1\le m\le 10^5\),
\(1\le k\le 10^5\).  Time limit 2 seconds, memory 256 MB.  One official sample:
the 10-vertex graph with edges
\(1\!-\!2,2\!-\!3,3\!-\!4,1\!-\!4,3\!-\!5,5\!-\!6,8\!-\!6,8\!-\!7,7\!-\!6,7\!-\!9,9\!-\!10\)
and queries \((1,2),(3,5),(6,9),(9,2),(9,3),(9,10)\) \(\to\) outputs
\(2,2,2,4,4,1\).
\end{problemstatement}

\begin{attemptgate}[Attempt it yourself first -- \lv{5}, give yourself 45 minutes]
Two vertices of a cactus can be joined by exponentially many simple paths -- each
cycle on the route doubles the count (two arcs around the ring).  So the answer is
\(2^{c}\) where \(c\) is the number of cycles you must cross between \(x\) and
\(y\).  The section's contraction theorem (\(\thmref{c31-cactustree}\)) turns
``count cycles on the \(x\)--\(y\) route'' into a tree question -- but which tree
object makes a whole cycle collapse to a single countable node, and what running
quantity answers \(c\) in \(\bigO{\log n}\) per query?\peek{Contract every maximal
cycle to one node; bridges stay as edges.  What remains is a tree on the
components (\(\thmref{c31-cactustree}\)).  Mark each contracted node ``heavy''
(it was a cycle).  Then \(c\) = number of heavy nodes on the tree path between the
component of \(x\) and the component of \(y\).}
\nudge{Find cycles by DFS: a back edge \(u\!-\!v\) (\(v\) an ancestor) plus the tree
path \(v\!\to\!u\) is one cycle -- union all its vertices into a DSU set and mark
them cyclic.  Vertex-disjointness means each vertex is unioned once, so it is
\(\bigO n\).  Build the component tree, root it, and store a root-to-node prefix
\(\mathrm{pre}[C]=\) (heavy nodes from root to \(C\) inclusive).}
\fullhint{Answer a query with LCA on the component tree:
\(c=\mathrm{pre}[C_x]+\mathrm{pre}[C_y]-2\,\mathrm{pre}[L]+\mathit{heavy}[L]\) where
\(L=\mathrm{lca}(C_x,C_y)\); print \(2^{c}\bmod 10^{9}{+}7\) (precompute powers of
two).  The \(+\mathit{heavy}[L]\) restores the LCA node subtracted twice.  Full
build and the six-query replay are below.}
\end{attemptgate}

\solstep{Understanding}
On a tree there is exactly one simple path between two vertices, so the count is
\(1=2^{0}\).  Every cycle is the only source of choice: entering a cycle at one
vertex and leaving at another, you may go around either side -- two edge-disjoint
arcs, a factor of \(2\).  Because a vertex cactus has \emph{vertex-disjoint}
cycles, the choices on different cycles are independent and each cycle contributes
exactly one factor of \(2\), no matter its length: the answer is \(2^{c}\) where
\(c\) is the number of cycles met along the \(x\)--\(y\) corridor.  Contracting
each cycle to one node (\(\thmref{c31-cactustree}\)) makes the corridor an ordinary
tree path, and ``how many cycle-nodes on the path'' is a textbook LCA + root-prefix
query (Ch.~16, \(\chref{c11}\)).

\solstep{Key Observations}
\begin{observation}
A cycle is a factor \(2\): between the entry vertex \(a\) and exit vertex \(b\) of a
cycle there are exactly two edge-simple arcs (clockwise, counter-clockwise), so any
route through that cycle doubles.  Routes are distinguished by \emph{edge sets}, so
the two arcs are always counted separately.
\end{observation}
\begin{observation}
Vertex-disjoint cycles \(\Rightarrow\) independence: crossing cycle \(A\) never uses an
edge of cycle \(B\), so choices multiply and \(c\) is additive -- the exponent counts
cycles, no double-counting worries.
\end{observation}
\begin{observation}
After contraction the graph is a tree, and a cycle-node is ``heavy''.  Making a root
prefix \(\mathrm{pre}[C]\) count heavy nodes on the root--\(C\) path turns
``heavy nodes on \(C_x\)--\(C_y\)'' into \(\mathrm{pre}[C_x]+\mathrm{pre}[C_y]-2\mathrm{pre}[L]+\mathit{heavy}[L]\)
with \(L=\mathrm{lca}(C_x,C_y)\) -- the standard vertex-path fold.
\end{observation}

\solstep{Think It With Me}
\begin{thinkbox}
\deadend{Enumerate the paths per query}{There can be \(2^{\Theta(n)}\) of them --
\emph{hopeless}.  You must count the exponent, not list the paths.}
\deadend{DFS the original graph for each query}{\(\bigO{(n+m)k}=\bigO{10^{10}}\) --
\emph{TLE}; and marking ``is this vertex on a cycle used here?'' per query is a mess.}
\deadend{Count bridges crossed}{Bridges are forced (one way) and carry no factor;
counting them gives the wrong exponent -- the choices come from \emph{cycles}, not
bridges.}
\aha{Contract each cycle to a heavy node (\(\thmref{c31-cactustree}\)); the answer is
\(2^{c}\) with \(c\) the heavy-node count on a \emph{tree} path -- one LCA and one
root-prefix, \(\bigO{\log n}\) per query after \(\bigO{n\log n}\) build.}
\end{thinkbox}

\solstep{The Approach Ladder}
\begin{center}\renewcommand{\arraystretch}{1.25}\setlength{\tabcolsep}{3pt}
\begin{tabular}{@{}c p{0.28\linewidth} p{0.30\linewidth} l c@{}}
\toprule
\textbf{\#} & \textbf{Approach} & \textbf{Idea} & \textbf{Time} & \textbf{Verdict}\\
\midrule
1 & \lv{2}\ list all paths & DFS, collect edge sets & exponential & \tle\\
2 & \lv{3}\ per-query graph DFS & walk \(x\to y\), count cycles & \(\bigO{(n+m)k}\) & \tle\\
3 & \lv{4}\ contract, \(O(n)\) path walk & DSU cycles + brute tree path & \(\bigO{nk}\) & \tle\\
4 & \lv{5}\ \textbf{contract + LCA prefix} & \(2^{\text{heavy-on-path}}\) & \(\bigO{n\log n+k\log n}\) & \ac\\
\bottomrule
\end{tabular}\end{center}
Rung 4 is the only one that survives \(n,m,k\sim 10^{5}\): pay \(\log n\) once per
query via binary lifting instead of walking the tree.

\solstep{Solution}
One DFS finds every cycle: a back edge to ancestor \(v\) makes the tree path
\(v\!\to\!u\) a cycle, so DSU-union its vertices and flag them cyclic
(vertex-disjointness bounds the total union work by \(\bigO n\)).  Collapse each DSU
set to a component node; bridges connect components and form the contracted tree
\(T\).  Root \(T\), mark a component heavy if it came from a cycle, build binary
lifting and the inclusive cycle-prefix \(\mathrm{pre}[C]\).  Per query, map
\(x,y\) to components \(C_x,C_y\), take \(L=\mathrm{lca}\), read
\(c=\mathrm{pre}[C_x]+\mathrm{pre}[C_y]-2\mathrm{pre}[L]+\mathit{heavy}[L]\), and
print \(2^{c}\bmod(10^{9}+7)\).

\solstep{Implementation}
\begin{lstlisting}[caption={CF 231E. DFS finds each cycle from a back edge, DSU-unions its vertices (total O(n) since cycles are vertex-disjoint); contract to a tree, binary-lift for LCA and fold a root-to-node cycle prefix; answer = 2 to the (heavy nodes on path) mod 1e9+7. ASCII identifiers, no underscores}, label={lst:c31-cactuscount}]
#include <bits/stdc++.h>
using namespace std;
const int MAXN=100005, B=18;
const long long MOD=1000000007LL;
int n, m, dsu[MAXN], disc[MAXN], par[MAXN], tm=0;
bool cycV[MAXN];
vector<int> adj[MAXN], T[MAXN];
int findp(int x){ return dsu[x]==x ? x : dsu[x]=findp(dsu[x]); }
void uni(int a,int b){ a=findp(a); b=findp(b); if(a!=b) dsu[b]=a; }
void dfsCactus(int u,int p){
  disc[u]=++tm; par[u]=p;
  for(size_t j=0;j<adj[u].size();j++){ int v=adj[u][j];
    if(v==p) continue;
    if(!disc[v]) dfsCactus(v,u);
    else if(disc[v]<disc[u]){              // back edge u-v: tree path v..u is a cycle
      int cur=u;
      while(cur!=v){ uni(cur,par[cur]); cycV[cur]=true; cur=par[cur]; }
      cycV[v]=true;                         // close onto the ancestor
    }
  }
}
int comp[MAXN], isCyc[MAXN], up[B][MAXN], depth[MAXN];
long long pre[MAXN], pw[MAXN];
int lca(int a,int b){
  if(depth[a]<depth[b]){ int t=a; a=b; b=t; }
  int d=depth[a]-depth[b];
  for(int j=0;j<B;j++) if((d>>j)&1) a=up[j][a];
  if(a==b) return a;
  for(int j=B-1;j>=0;j--) if(up[j][a]!=up[j][b]){ a=up[j][a]; b=up[j][b]; }
  return up[0][a];
}
int main(){
  scanf("%d %d", &n, &m);
  for(int v=1;v<=n;v++) dsu[v]=v;
  for(int i=0;i<m;i++){ int a,b; scanf("%d %d", &a, &b);
    adj[a].push_back(b); adj[b].push_back(a); }
  dfsCactus(1,0);
  for(int v=1;v<=n;v++){ comp[v]=findp(v); if(cycV[v]) isCyc[comp[v]]=true; }
  for(int u=1;u<=n;u++) for(size_t j=0;j<adj[u].size();j++){ int v=adj[u][j];
    if(u<v){ int cu=comp[u], cv=comp[v];
      if(cu!=cv){ T[cu].push_back(cv); T[cv].push_back(cu); } } }
  // iterative DFS over the contracted tree from the component of vertex 1
  int root=comp[1]; vector<char> seen(n+1,0); vector<int> order;
  seen[root]=1; up[0][root]=0; depth[root]=0; pre[root]= isCyc[root]?1:0;
  order.push_back(root);
  for(size_t i=0;i<order.size();i++){ int c=order[i];
    for(size_t j=0;j<T[c].size();j++){ int d=T[c][j]; if(seen[d]) continue; seen[d]=1;
      up[0][d]=c; depth[d]=depth[c]+1; pre[d]=pre[c]+(isCyc[d]?1:0);
      order.push_back(d); } }
  for(int j=1;j<B;j++) for(int c=1;c<=n;c++) up[j][c]=up[j-1][ up[j-1][c] ];
  pw[0]=1; for(int i=1;i<=n;i++) pw[i]=pw[i-1]*2%MOD;
  int k; scanf("%d", &k);
  while(k--){ int x,y; scanf("%d %d", &x, &y);
    int cx=comp[x], cy=comp[y], L=lca(cx,cy);
    long long c=pre[cx]+pre[cy]-2*pre[L]+(isCyc[L]?1:0);
    printf("%lld\n", pw[c]);
  }
  return 0;
}
\end{lstlisting}

\solstep{Trace on the Official Samples}
\begin{tracebox}[Contracting the sample graph, then all six queries]
The two cycles are \(A=\{1,2,3,4\}\) (the back edge \(1\!-\!4\) over tree path \(4\!-\!3\!-\!2\!-\!1\))
and \(B=\{6,7,8\}\); vertices \(5,9,10\) are singletons.  Contracted tree,
root \(A\): \(A - 5 - B - 9 - 10\), with heavy nodes \(A\) and \(B\).
Cycle-prefix \(\mathrm{pre}\): \(\mathrm{pre}[A]{=}1,\ \mathrm{pre}[5]{=}1,\ \mathrm{pre}[B]{=}2,\ \mathrm{pre}[9]{=}2,\ \mathrm{pre}[10]{=}2\).\quad
\textbf{\(c=\mathrm{pre}[C_x]+\mathrm{pre}[C_y]-2\mathrm{pre}[L]+\mathit{heavy}[L]\)}:
\begin{itemize}[nosep,leftmargin=1.4em]
  \item \((1,2)\): both \(\to A\), \(L{=}A\): \(1{+}1{-}2{+}1=1\) \(\to 2^{1}=\mathbf{2}\).
  \item \((3,5)\): \(A,5\), \(L{=}A\): \(1{+}1{-}2{+}1=1\) \(\to \mathbf{2}\).
  \item \((6,9)\): \(B,9\), \(L{=}B\): \(2{+}2{-}4{+}1=1\) \(\to \mathbf{2}\).
  \item \((9,2)\): \(9,A\), \(L{=}A\): \(2{+}1{-}2{+}1=2\) \(\to 2^{2}=\mathbf{4}\).
  \item \((9,3)\): \(9,A\), \(L{=}A\): same \(=2\) \(\to \mathbf{4}\).
  \item \((9,10)\): \(9,10\), \(L{=}9\) (not heavy): \(2{+}2{-}4{+}0=0\) \(\to \mathbf{1}\).
\end{itemize}
Outputs \(2,2,2,4,4,1\) -- exactly the official sample. \(\Box\)
\end{tracebox}

\solstep{Complexity}
\begin{complexitybox}
Cycle DFS is \(\bigO{n+m}\): each vertex is unioned at most once because the cycles
are vertex-disjoint.  Building \(T\) is \(\bigO{n+m}\); the lifting table is
\(\bigO{n\log n}\) and a query (one LCA) is \(\bigO{\log n}\).  Total
\(\bigO{n\log n + k\log n}\) \(\approx 3.4\cdot 10^{6}\) for \(n=m=k=10^{5}\), inside 2
seconds.  Memory is the \(B\times n\) lift table \(\approx 1.8\cdot 10^{6}\) ints
\(\approx 7\) MB, well under 256 MB.  The exponent \(c\le\) (number of cycles) \(\le n\),
so \(\mathrm{pw}\) is precomputed to \(n\) modulo \(10^{9}{+}7\).
\end{complexitybox}

\begin{hardeningbox}
\begin{itemize}[nosep,leftmargin=1.4em]
\item \textbf{Add the LCA node back.}  The prefix counts heavy nodes \emph{inclusive},
        so \(-2\,\mathrm{pre}[L]\) removes \(L\) twice; you must add
        \(\mathit{heavy}[L]\) back.  Forgetting it under-counts whenever the lowest
        common component is itself a cycle.
\item \textbf{Recursion depth on the DFS.}  A cactus can be a long path; the
        \code{dfsCactus} recursion may near \(10^{5}\) frames.  Raise the stack, or
        convert to the iterative form used for \(T\) in \(\coderef{c31-cactuscount}\).
\item \textbf{Back-edge detection needs \(v\ne\) parent.}  In a simple graph
        \code{if(v==p) continue} is enough; there are no parallel edges to confuse the
        tree edge with a 2-cycle, so every non-parent visited neighbour with
        \(\mathrm{disc}[v]<\mathrm{disc}[u]\) is a genuine back edge.
\item \textbf{Powers of two modulo.}  Precompute \(\mathrm{pw}[i]=2^{i}\bmod(10^{9}{+}7)\)
        once; recomputing per query with \texttt{pow}/bit tricks invites overflow and
        precision bugs.
\end{itemize}
\end{hardeningbox}

\begin{recallbox}[29.4 in one breath]
\emph{A cactus is a graph whose blocks are only bridges and simple cycles -- every edge
on at most one cycle (\(\defref{c31-cactus}\)).  Contract each cycle to a node and the
block-cut tree is a \textbf{tree} (\(\thmref{c31-cactustree}\)), so you DP on it with one
new primitive: the cyclic-order transfer \code{ring} (\(\lemref{c31-cycleorder}\)) that
pins the entry vertex and folds the ring, closing only on a consistent last edge.  The
bowtie counts to \(10\) independent sets (\(\figref{c31-cactus}\)); swap the state table
for colouring or domination and the same skeleton answers those.}
\selftest{
  \item What blocks are allowed in a cactus? \emph{Bridges and single simple cycles only; a block that is a denser 2-connected blob (e.g. a shared-edge diamond) breaks the cactus property \(\defref{c31-cactus}\).)}
  \item Why does contracting cycles leave a tree? \emph{Cycles meet only at cut vertices (no shared edge), and the block-cut tree already has no cycles among block/cut nodes \(\thmref{c31-cactustree}\).)}
  \item What extra step does a cycle block add over a tree DP? \emph{A ring transfer that pins the entry vertex's state and folds the cycle in order, rejecting the wrap-around conflict \(\lemref{c31-cycleorder}\).)}
  \item In CF 231E, why is the answer a power of two and what is the exponent? \emph{Each cycle on the \(x\)--\(y\) route offers two edge-disjoint arcs, and vertex-disjoint cycles make the choices independent, so it is \(2^{c}\) with \(c\) the number of cycle-nodes on the contracted-tree path \(\probref{c31-cactuscount}\).)}
}
\end{recallbox}
